% sections/01_introduction.tex
% Reordered: result 1 (non-neutrality) is the headline; results 2 (distribution)
% and 3 (ex-ante) are subordinated; EU framing committed; booking + supply-elasticity
% caveats collected in one paragraph (referee m5); sigma_eps discipline stated once.
%
% Expected labels (align to your files): sec:literature sec:model sec:calibration
% sec:channel sec:shielding sec:exante sec:policy sec:conclusion app:booking app:taste
% Bibitem keys: Bruegel2023 EHPA2023 EHPA2024 EHPA2025 auclert2023energy langot2026tariff bayer2026hank
%               GallagherMuehlegger2011 Chandra2010 MuehleggerRapson2022

\section{Introduction}\label{sec:intro}

When wholesale energy prices tripled in 2021--22, European governments responded
by shielding consumers. Between September 2021 and January 2023 EU member states
committed more than \euro{}600 billion to protect households and firms, mostly
through measures that hold down the retail price itself: caps, regulated tariffs,
and energy tax cuts \citep{Bruegel2023}. These measures protected purchasing
power, which was their purpose. But they act on the price that governs the energy
transition. A retail price that is held fixed cannot signal scarcity, and
households substitute away from fossil energy in response to that signal.
Bruegel, which assembled the spending figures, drew the same conclusion and urged
governments to replace price-suppressing measures, de facto fossil-fuel
subsidies, with targeted income support \citep{Bruegel2023}.

This paper studies one margin of that trade-off which the macroeconomic
literature has left out: the adoption of green durables during the crisis. A
fossil price spike is the moment when switching to a technology that
does not burn fossil energy pays best, and European households responded along
this margin. Heat-pump sales reached a record of about three million units in
2022, up roughly 38\% as gas prices surged. They fell 5\% in 2023 and a further
21\% in 2024 as gas prices normalised; the industry association named the low
price of subsidised gas among the three main causes, and sales recovered only
once governments cut electricity taxes \citep{EHPA2023,EHPA2024,EHPA2025}. The
margin responds to the relative price, and the price shields reached it.

We introduce this margin into an otherwise standard heterogeneous-agent model of
an energy-importing economy. Households in the small open economy HANK of
Auclert, Monnery, Rognlie, and Straub (\citeyear{auclert2023energy}; henceforth AMRS)
choose discretely whether to pay a one-time resource cost for a green durable, a
heat pump or an electric vehicle, that insulates them from the fossil price. The
choice is a logit over the durable stock and is solved in the sequence space
\citep{Auclert2021a}.
Everything else is inherited from \citet{auclert2023energy}: the CES energy nest, sticky
import prices, the wage Phillips curve, and the constant-real-rate monetary
baseline. The aggregate block is a generic energy importer; the extensive margin
is the only new object, and we calibrate to the European episode
(Section~\ref{sec:cal}, Section~\ref{sec:policy}). The closest
antecedent is \citet{dietrich2025consumer}, who shows in a representative-agent model
that a carbon-pricing transition is mildly inflationary and that green durables
are interest-rate sensitive. We keep that channel but add household
heterogeneity, an energy shock, and the fiscal and open-economy blocks. Our
question is not how a transition unfolds but how a crisis, and the policies that
meet it, move the adoption margin.

The central result is a non-neutrality. An energy shock endogenously triggers an
adoption wave: in the baseline crisis the green share rises by about eight
percentage points at the peak, from a five-percent base. The two ex-post
instruments in the recent literature, a retail price cap and a
Slutsky-compensating transfer indexed to pre-crisis energy use, stabilise output
comparably but differ on the new margin. The cap holds the brown retail price at
its pre-crisis level and thereby switches adoption off: the peak green-share
response falls from eight points to essentially zero. The transfer supports
incomes and leaves the wave intact. Shielding consumers from the relative price
is therefore not neutral for the transition, and the cost is the loss of the
extensive margin itself rather than a marginal distortion of intensive-margin
demand. The mechanism is the brown retail price: the cap fixes it and the
transfer leaves it free, so the ranking survives every other modelling choice
(Section~\ref{sec:shield}).

The second question is who should be shielded, and how. Energy demand is close
to proportional to total consumption within technology groups. A transfer
indexed to pre-crisis energy use therefore hands the largest cheques to the
wealthiest households, which have the lowest MPCs and gain the least insurance per euro. An untargeted flat transfer of the same fiscal cost
reaches the high-MPC households the crisis hits hardest and delivers strictly
higher welfare (a consumption-equivalent variation of $-0.36\%$ against
$-0.84\%$), while leaving adoption intact. This result is conditional on
homothetic energy demand. With a declining Engel curve for energy, as in
\citet{bayer2026hank} through energy-share heterogeneity or in \citet{langot2026tariff}
through a subsistence floor, low-income households carry a larger energy share
and the ranking could weaken. We take the homothetic nest from \citet{auclert2023energy}
and leave that case aside. Under constant energy shares, energy-indexed
targeting is dominated by an equal-cost flat transfer.

The third question moves from ex-post to ex-ante policy. Europe will soon run
the ex-ante experiment deliberately: the forthcoming ETS2 carbon price for
buildings and road transport will raise the household brown-energy price by
design \citep{ECETS2}. A permanent carbon tax that greens the steady state
before the crisis is a poor substitute for ex-post insurance. Its crisis
dividend is zero to first order, and it does not improve as shocks become more
persistent. The prepared economy in fact adopts \emph{more} during the crisis; what falls is
the welfare value of the marginal switch, because pre-greening uses up, before
the shock, an adjustment margin that is most valuable when the economy enters
the crisis with room to green (Section~\ref{sec:exante}). The externality and
the standing gain of a greener steady state justify the carbon tax; insurance
against the next shock does not.

Two auxiliary results are conditional, and we state the conditions rather than
resolve them. The sign of the cap's private-welfare effect depends on the
energy-supply elasticity, which locates one axis of the disagreement between
\citet{langot2026tariff} and \citet{bayer2026hank} without nesting their models. The sign
of the adoption channel's contribution to output depends on how the green
resource flows are booked in the balance of payments; we report the import
booking throughout and collect the alternative in Appendix~\ref{app:booking}.
The policy ordering is robust to both choices, since it follows from the
relative price. The size of the adoption wave, by contrast, depends on one
parameter, the logit taste scale $\sigma_\varepsilon$, and we discipline it with
data rather than assume it. At a five-percent green share every brown household
sits in the exponential tail of the logit, so the semi-elasticity of switching
is $1/\sigma_\varepsilon$ household by household (Appendix~\ref{app:taste}).
Matching the response of adoption to a \$1{,}000 incentive in the
quasi-experimental vehicle literature \citep{GallagherMuehlegger2011,
Chandra2010, MuehleggerRapson2022} pins $\sigma_\varepsilon\in[0.05,0.07]$, with
an implied green premium of \$15{,}000--\$18{,}000. The steady state and the
policy ordering do not depend on the scale; the size of the wave does.

Section~\ref{sec:lit} places the paper in the literature.
Section~\ref{sec:model} presents the model and Section~\ref{sec:cal} the
calibration, the identification of the taste scale, and the experimental design (the
two energy shocks, the frozen-choice counterfactual that isolates the adoption
channel, the policy instruments, and the welfare measure). Section~\ref{sec:crisis} documents the
crisis adoption wave and its contribution to the aggregate response.
Section~\ref{sec:shield} compares the ex-post instruments, prices against incomes,
and Section~\ref{sec:exante} the ex-ante carbon tax. Section~\ref{sec:synthesis}
ranks the instruments by welfare and states what the results depend on.
Section~\ref{sec:policy} maps the results into European policy, and
Section~\ref{sec:conclusion} concludes.
